\section*{Capítulo \ref{ch:g5:raw-materials-and-recycling} --- Das matérias-primas à reciclagem}

\begin{solution}{exo:g5:raw-materials-and-recycling:1}
Naturais: madeira, lã, pedra. Fabricados: vidro, aço, plástico.
\end{solution}

\begin{solution}{exo:g5:raw-materials-and-recycling:2}
Papel: madeira (árvores). Vidro: areia (com barrilha e calcário).
Algodão: o algodoeiro. Plástico: petróleo.
\end{solution}

\begin{solution}{exo:g5:raw-materials-and-recycling:3}
Uma rocha da qual se extrai um metal. O alumínio vem da bauxita.
\end{solution}

\begin{solution}{exo:g5:raw-materials-and-recycling:4}
O pote no coletor de vidro, o jornal com o papel e o papelão, a latinha
com os metais, a garrafa com as garrafas de plástico.
\end{solution}

\begin{solution}{exo:g5:raw-materials-and-recycling:5}
A \omterm{def:g5:raw-materials-and-recycling:raw-material}{matéria-prima} nova: o material reciclado é usado no lugar da areia, do
\omterm{def:g5:raw-materials-and-recycling:raw-material}{minério} ou do petróleo tirados da natureza.
\end{solution}

\begin{solution}{exo:g5:raw-materials-and-recycling:6}
Renováveis: madeira (se replantada), lã, algodão. Não renováveis:
petróleo, \omterm{def:g5:raw-materials-and-recycling:raw-material}{minério} de ferro.
\end{solution}

\begin{solution}{exo:g5:raw-materials-and-recycling:7}
Produzir vidro a partir de areia é uma \omterm{def:g5:irreversible-changes:chemical-change}{transformação química}: algo novo
é fabricado, e não há um caminho simples de volta.
\end{solution}

\begin{solution}{exo:g5:raw-materials-and-recycling:8}
Por exemplo: reduzir --- comprar frutas sem embalagem de plástico;
reutilizar --- encher de novo um pote de vidro; reciclar --- pôr uma
latinha no coletor de metais.
\end{solution}

\begin{solution}{exo:g5:raw-materials-and-recycling:9}
Cada material é reciclado do seu próprio jeito: pedaços de outro
material estragariam o vidro derretido.
\end{solution}

\begin{solution}{exo:g5:raw-materials-and-recycling:10}
Um vinte avos de 20 unidades: $20 \div 20 = 1$ unidade de energia, mais ou menos.
\end{solution}

\begin{solution}{exo:g5:raw-materials-and-recycling:11}
A \omterm{def:g5:raw-materials-and-recycling:recycling}{reciclagem} ainda precisa de energia, de máquinas e de transporte; um
objeto que nunca é produzido poupa a \omterm{def:g5:raw-materials-and-recycling:raw-material}{matéria-prima} dele e toda essa
energia.
\end{solution}
